% Reviewed translation: PDF pages 219--227 / printed pages 205--213.
% The original scan is authoritative; mathematical tokens were checked against it.
\MathBookSourcePage{219}
\subsection{混合逼近及其在 $l$ 次有限元上的应用}
\label{subsec:incompressible-sfvp-mixed-approximation}

在上一小节的记号下, 置
\[
\widetilde X
=\{\mathbf v=(\operatorname{curl}\phi,\theta);
\ \phi\in\Phi_s,\ \theta\in L^2(\Omega)\},
\qquad
\widetilde M=W^{1,r}(\Omega),
\]
并赋予范数
\[
\| (\operatorname{curl}\phi,\theta)\|_{\widetilde X}
=|\phi|_{1,s,\Omega}+\|\theta\|_{0,\Omega},
\qquad
\|\mu\|_{\widetilde M}=\|\mu\|_{1,r,\Omega}.
\]
双线性形式为
\[
\widetilde a(\mathbf u,\mathbf v)=\nu(\omega,\theta)
\quad
\forall\mathbf u=(\operatorname{curl}\psi,\omega),
\ \mathbf v=(\operatorname{curl}\phi,\theta)\in\widetilde X,
\]
以及
\[
\widetilde b(\mathbf v,\mu)
=(\operatorname{curl}\phi,\operatorname{curl}\mu)-(\theta,\mu)
\quad
\forall\mathbf v=(\operatorname{curl}\phi,\theta)\in\widetilde X,
\ \forall\mu\in\widetilde M.
\]
相应的空间为
\[
\widetilde V
=\{\mathbf v\in\widetilde X;
\ \widetilde b(\mathbf v,\mu)=0
\ \forall\mu\in\widetilde M\}.
\]

为构造逼近, 引入三个有限维空间
\begin{equation}\label{eq:incompressible-sfvp-discrete-spaces}
\Phi_h\subset\Phi_s,
\qquad
\Theta_h\subset L^2(\Omega),
\qquad
M_h\subset W^{1,r}(\Omega),
\end{equation}
并把问题 \eqref{eq:incompressible-weak-stream-three-field}
离散化为下列问题 $(\widetilde Q_h)$:

求 $(\psi_h,\omega_h)\in\Phi_h\times\Theta_h$
和 $\lambda_h\in M_h$, 使
\begin{equation}\label{eq:incompressible-sfvp-discrete-four-field}
\left\{
\begin{aligned}
\nu(\omega_h,\theta_h)-(\lambda_h,\theta_h)
+(\operatorname{curl}\phi_h,\operatorname{curl}\lambda_h)
&=(\mathbf f,\operatorname{curl}\phi_h)
&&\forall(\phi_h,\theta_h)\in\Phi_h\times\Theta_h,\\
(\operatorname{curl}\psi_h,\operatorname{curl}\mu_h)
&=(\omega_h,\mu_h)
&&\forall\mu_h\in M_h.
\end{aligned}
\right.
\end{equation}
这里
\[
\widetilde X_h
=\{\mathbf v_h=(\operatorname{curl}\phi_h,\theta_h);
\ \phi_h\in\Phi_h,\ \theta_h\in\Theta_h\},
\]
而
\[
\widetilde V_h
=\{\mathbf v_h=(\operatorname{curl}\phi_h,\theta_h)
\in\widetilde X_h;
\ (\operatorname{curl}\phi_h,\operatorname{curl}\mu_h)
=(\theta_h,\mu_h)\ \forall\mu_h\in M_h\}.
\]
下一个 Lemma 验证 inf--sup 条件
\eqref{eq:incompressible-discrete-inf-sup}.

\begin{Lemma}\label{lem:incompressible-sfvp-discrete-inf-sup}
若 $M_h\subset\Theta_h$, 则
\begin{equation}\label{eq:incompressible-sfvp-discrete-inf-sup}
\sup_{\mathbf v_h\in\widetilde X_h}
\frac{\widetilde b(\mathbf v_h,\mu_h)}
{\|\mathbf v_h\|_{\widetilde X}}
\geq\|\mu_h\|_{0,\Omega}
\qquad\forall\mu_h\in M_h.
\end{equation}
\end{Lemma}

事实上, \eqref{eq:incompressible-sfvp-discrete-inf-sup}
由取 $\mathbf v_h=(0,-\mu_h)$ 得到.

当 $s=2$ 时, $\widetilde a(\cdot,\cdot)$ 的一致椭圆性
由下一个 Lemma 推出. 但当 $s>2$ 时, 似乎没有明显的方法证明这一性质,
因而必须把它作为一个假设引入.

\begin{Lemma}\label{lem:incompressible-sfvp-discrete-stream-bound}
若 $\Phi_h\subset M_h$, 则存在常数 $C>0$, 使
\begin{equation}\label{eq:incompressible-sfvp-discrete-stream-bound}
|\phi_h|_{1,\Omega}
\leq C\|\theta_h\|_{0,\Omega}
\qquad
\forall(\operatorname{curl}\phi_h,\theta_h)\in\widetilde V_h.
\end{equation}
\end{Lemma}

\begin{Proof}
由定义, $\widetilde V_h$ 中的函数
$(\operatorname{curl}\phi_h,\theta_h)$ 满足
\[
(\operatorname{curl}\phi_h,\operatorname{curl}\mu_h)
=(\theta_h,\mu_h)
\qquad\forall\mu_h\in M_h.
\]
由于 $\Phi_h\subset M_h$, 可以取 $\mu_h=\phi_h$.
再利用 Lemma~\ref{lem:sec3-boundary-poincare}, 得
\[
|\phi_h|_{1,\Omega}^2
\leq\|\theta_h\|_{0,\Omega}\|\phi_h\|_{0,\Omega}
\leq C\|\phi_h\|_{1,\Omega}\|\theta_h\|_{0,\Omega},
\]
从而得到 \eqref{eq:incompressible-sfvp-discrete-stream-bound}.
\end{Proof}

\MathBookSourcePage{220}
当 $s=2$ 时, 由该 Lemma 可知, 只要 $\Phi_h\subset M_h$,
就存在 $\alpha^*>0$, 使
\begin{equation}\label{eq:incompressible-sfvp-discrete-ellipticity}
\nu\|\theta_h\|_{0,\Omega}^2
\geq\alpha^*\|\mathbf v_h\|_{\widetilde X}^2
\qquad
\forall\mathbf v_h=(\operatorname{curl}\phi_h,\theta_h)
\in\widetilde V_h.
\end{equation}
换言之, \eqref{eq:incompressible-vorticity-seminorm} 定义的半范数
$|\cdot|$ 与 $\|\cdot\|_{\widetilde X}$
在 $\widetilde V_h$ 上一致等价.
当 $s>2$ 时, 我们把这一性质作为假设, 并将在后面的例子中验证它.

\medskip
\noindent\textbf{Hypothesis H1.}
当 $s>2$ 时, $|\cdot|$ 与 $\|\cdot\|_{\widetilde X}$
是在 $\widetilde V_h$ 上关于 $h$ 一致等价的两个范数.

\begin{Lemma}\label{lem:incompressible-sfvp-discrete-wellposed}
\textup{1)} 假设
\begin{equation}\label{eq:incompressible-sfvp-space-inclusions}
\Phi_h\subset M_h\subset\Theta_h,
\end{equation}
并且当 $s>2$ 时假设 Hypothesis H1 成立.
则问题 \eqref{eq:incompressible-sfvp-discrete-four-field}
有唯一解
$\mathbf u_h=(\operatorname{curl}\psi_h,\omega_h)
\in\widetilde X_h$ 和 $\lambda_h\in M_h$.

\textup{2)} 若此外 $M_h=\Theta_h$, 则 $\lambda_h=\nu\omega_h$.
\end{Lemma}

\begin{Proof}
\textup{1)} 由 Lemma~\ref{lem:incompressible-sfvp-discrete-inf-sup},
Lemma~\ref{lem:incompressible-sfvp-discrete-stream-bound}, Hypothesis H1
和 Theorem~\ref{thm:incompressible-abstract-mixed-error} 立即得到第一部分.

\textup{2)} 令 $M_h=\Theta_h$.
我们证明 $(\mathbf u_h,\nu\omega_h)$
满足 \eqref{eq:incompressible-sfvp-discrete-four-field}, 即
\[
\nu(\operatorname{curl}\omega_h,\operatorname{curl}\phi_h)
=(\mathbf f,\operatorname{curl}\phi_h)
\qquad\forall\phi_h\in\Phi_h.
\]
若 $(\operatorname{curl}\phi_h,\theta_h)\in\widetilde V_h$,
则 \eqref{eq:incompressible-sfvp-discrete-four-field} 化为
\[
\nu(\omega_h,\theta_h)
=(\mathbf f,\operatorname{curl}\phi_h).
\]
但由于 $M_h=\Theta_h$, $\phi_h$ 和 $\theta_h$ 特别满足
\[
(\operatorname{curl}\phi_h,\operatorname{curl}\omega_h)
=(\theta_h,\omega_h).
\]
因此
\[
\nu(\operatorname{curl}\phi_h,\operatorname{curl}\omega_h)
=(\mathbf f,\operatorname{curl}\phi_h)
\quad
\forall(\operatorname{curl}\phi_h,\theta_h)\in\widetilde V_h.
\]
这正是所需结论, 因为当 $M_h=\Theta_h$ 时,
每个 $\phi_h\in\Phi_h$ 恰有一个 $\theta_h\in\Theta_h$,
使 $(\operatorname{curl}\phi_h,\theta_h)\in\widetilde V_h$.
\end{Proof}

注意, 当 $s>2$ 时, 上述证明不要求 $\Phi_h\subset M_h$,
但实际中这一包含关系总是成立.

Lemma~\ref{lem:incompressible-sfvp-discrete-wellposed}
有下列重要推论.

\begin{Theorem}\label{thm:incompressible-sfvp-practical-equivalence}
在 Lemma~\ref{lem:incompressible-sfvp-discrete-wellposed} 的假设下,
问题 \eqref{eq:incompressible-sfvp-discrete-four-field} 等价于:

求 $\psi_h\in\Phi_h$ 和 $\omega_h\in\Theta_h$, 使
\begin{equation}\label{eq:incompressible-sfvp-practical-problem}
\left\{
\begin{aligned}
\nu(\operatorname{curl}\omega_h,\operatorname{curl}\phi_h)
&=(\mathbf f,\operatorname{curl}\phi_h)
&&\forall\phi_h\in\Phi_h,\\
(\operatorname{curl}\psi_h,\operatorname{curl}\mu_h)
&=(\omega_h,\mu_h)
&&\forall\mu_h\in\Theta_h.
\end{aligned}
\right.
\end{equation}
\end{Theorem}

\MathBookSourcePage{221}
问题 \eqref{eq:incompressible-sfvp-practical-problem}
是实际计算所依据的表述. 它与连续情形一样具有下列吸引人的特点:

\textup{1)} 问题 \eqref{eq:incompressible-sfvp-practical-problem}
以具有 $H^1$ 正则性的流函数来表述, 因而速度严格无散;

\textup{2)} 它由两个关于 Laplace 算子的 Dirichlet 问题组成;

\textup{3)} 流函数和涡量的计算独立于以后任何压力计算.

同样, \eqref{eq:incompressible-sfvp-practical-problem}
中的两个问题是耦合的. 它们既可以一起求解,
也可以用与 Subsections~\ref{subsec:primitive-glowinski-pironneau}
和 \ref{subsec:primitive-gp-implementation} 中所讨论技术非常相似的方法分离求解
(参见 Glowinski \& Pironneau~\cite{ref:glowinski-pironneau-39}).

由 Theorem~\ref{thm:incompressible-abstract-mixed-error},
条件 \eqref{eq:incompressible-sfvp-space-inclusions} 与 Hypothesis H1
蕴含误差估计
\eqref{eq:incompressible-velocity-abstract-error} 和
\eqref{eq:incompressible-multiplier-abstract-error},
其中的常数与 $h$ 无关.
此外, 若 $M_h=\Theta_h$,
则可引入由 \eqref{eq:appendix-neumann-ritz-projection}
定义的从 $W^{1,r}(\Omega)$ 到 $\Theta_h$ 的投影算子 $P_h$,
从而略微改进 \eqref{eq:incompressible-multiplier-abstract-error} 的界:
\[
(\operatorname{grad}(P_h\mu-\mu),\operatorname{grad}\theta_h)=0
\quad\forall\theta_h\in\Theta_h,
\qquad
(P_h\mu-\mu,1)=0.
\]
所得估计由下一个 Theorem 给出.

\begin{Theorem}\label{thm:incompressible-sfvp-velocity-vorticity-error}
令
$\mathbf u=(\operatorname{curl}\psi,\omega=-\Delta\psi)$
和 $\mathbf u_h=(\operatorname{curl}\psi_h,\omega_h)$
分别为问题 \eqref{eq:incompressible-stream-vorticity-pair}
和 \eqref{eq:incompressible-sfvp-practical-problem} 的解.
在 Corollary~\ref{cor:incompressible-weak-stream-unique}
和 Lemma~\ref{lem:incompressible-sfvp-discrete-wellposed} 的假设下,
有误差界
\begin{equation}\label{eq:incompressible-sfvp-quasioptimal-error}
\|\mathbf u-\mathbf u_h\|_{\widetilde X}
\leq C_1\left\{
\inf_{\mathbf v_h\in\widetilde V_h}
\|\mathbf u-\mathbf v_h\|_{\widetilde X}
+\|\omega-P_h\omega\|_{0,\Omega}
\right\},
\end{equation}
以及
\begin{equation}\label{eq:incompressible-sfvp-expanded-error}
\begin{split}
\|\mathbf u-\mathbf u_h\|_{\widetilde X}
\leq C_2\biggl[&(1+K_r(h))
\left\{
\inf_{\phi_h\in\Phi_h}|\psi-\phi_h|_{1,s,\Omega}
+\inf_{\theta_h\in\Theta_h}
\|\omega-\theta_h\|_{0,\Omega}
\right\}\\
&+\|\omega-P_h\omega\|_{0,\Omega}\biggr],
\end{split}
\end{equation}
其中 $P_h$ 由 \eqref{eq:appendix-neumann-ritz-projection} 定义,
而对 $r<2$,
\[
K_r(h)=\sup_{\theta_h\in\Theta_h}
\frac{\|\theta_h\|_{1,r,\Omega}}
{\|\theta_h\|_{0,\Omega}}.
\]
\end{Theorem}

\begin{Proof}
先证明 \eqref{eq:incompressible-sfvp-quasioptimal-error}.
由 \eqref{eq:incompressible-stream-vorticity-pair}
和 \eqref{eq:incompressible-sfvp-practical-problem} 得
\[
(\operatorname{curl}(\psi-\psi_h),\operatorname{curl}\phi_h)
=(\omega-\omega_h,\mu_h)
\quad\forall\phi_h\in\Phi_h,
\ \forall\mu_h\in\Theta_h.
\]
特别地, 取 $\mu_h=P_h\omega$, 得
\begin{equation}\label{eq:incompressible-sfvp-projected-orthogonality}
(\operatorname{curl}(P_h\omega-\omega_h),
\operatorname{curl}\phi_h)=0
\qquad\forall\phi_h\in\Phi_h.
\end{equation}

\MathBookSourcePage{222}
因此, $\widetilde V_h$ 的定义蕴含
\begin{equation}\label{eq:incompressible-sfvp-vorticity-orthogonality}
(P_h\omega-\omega_h,\theta_h)=0
\qquad
\forall\mathbf v_h=(\operatorname{curl}\phi_h,\theta_h)
\in\widetilde V_h.
\end{equation}
从而
\begin{equation}\label{eq:incompressible-sfvp-vorticity-best-approximation}
\|\omega_h-\theta_h\|_{0,\Omega}
\leq\|P_h\omega-\theta_h\|_{0,\Omega}
\qquad
\forall\mathbf v_h=(\operatorname{curl}\phi_h,\theta_h)
\in\widetilde V_h.
\end{equation}
$|\cdot|$ 与 $\|\cdot\|_{\widetilde X}$
在 $\widetilde V_h$ 上的等价性于是给出
\[
\|\mathbf u_h-\mathbf v_h\|_{\widetilde X}
\leq C\left\{
\|P_h\omega-\omega\|_{0,\Omega}
+\|\mathbf u-\mathbf v_h\|_{\widetilde X}
\right\}
\qquad\forall\mathbf v_h\in\widetilde V_h.
\]
这证明了 \eqref{eq:incompressible-sfvp-quasioptimal-error}.
估计 \eqref{eq:incompressible-sfvp-expanded-error}
由 \eqref{eq:incompressible-sfvp-quasioptimal-error}
和 \eqref{eq:incompressible-affine-space-approximation} 立即得到.
\end{Proof}

\begin{Remark}\label{rem:incompressible-sfvp-sharper-seminorm-error}
公式 \eqref{eq:incompressible-sfvp-vorticity-best-approximation}
给出 $|\mathbf u-\mathbf u_h|$ 的略尖锐估计
\[
|\mathbf u-\mathbf u_h|
\leq2\inf_{\mathbf v_h\in\widetilde V_h}
|\mathbf u-\mathbf v_h|
+\|\omega-P_h\omega\|_{0,\Omega}.
\]
\end{Remark}

现在转向压力. 为求解压力, 引入第四个有限维空间
\begin{equation}\label{eq:incompressible-sfvp-pressure-space}
Q_h\subset W^{1,s}(\Omega)\cap L_0^2(\Omega),
\end{equation}
并据此直接离散问题 \eqref{eq:incompressible-decoupled-pressure}:

求 $p_h\in Q_h$, 使
\begin{equation}\label{eq:incompressible-sfvp-discrete-pressure}
(\operatorname{grad}p_h,\operatorname{grad}q_h)
=(\mathbf f-\nu\operatorname{curl}\omega_h,
\operatorname{grad}q_h)
\qquad\forall q_h\in Q_h.
\end{equation}
对固定的 $\omega_h$, 这个方阵有唯一解,
因为 $\operatorname{grad}p=0$ 且 $p\in L_0^2(\Omega)$ 蕴含 $p=0$.

\begin{Remark}\label{rem:incompressible-sfvp-combined-problem}
问题 \eqref{eq:incompressible-sfvp-practical-problem}
和 \eqref{eq:incompressible-sfvp-discrete-pressure}
可以合并成一个问题:

求 $\psi_h\in\Phi_h$, $\omega_h\in\Theta_h$ 和 $p_h\in Q_h$, 使
\[
\nu(\omega_h,\operatorname{curl}\mathbf v_h)
+(\operatorname{grad}p_h,\mathbf v_h)
=(\mathbf f,\mathbf v_h)
\quad
\forall\mathbf v_h=\operatorname{curl}\phi_h+\operatorname{grad}q_h,
\quad \phi_h\in\Phi_h,\ q_h\in Q_h,
\]
以及
\[
(\operatorname{curl}\psi_h,\operatorname{curl}\mu_h)
=(\omega_h,\mu_h)
\qquad\forall\mu_h\in\Theta_h.
\]
注意它与问题 \eqref{eq:incompressible-hhj-velocity-pressure-weak}
之间的类比.
\end{Remark}

\begin{Remark}\label{rem:incompressible-sfvp-discrete-neumann}
与连续情形相同, 问题
\eqref{eq:incompressible-sfvp-discrete-pressure}
使用具有 $H^1$ 正则性的压力, 并求解 Laplace 算子的一个离散 Neumann 问题.
\end{Remark}

为估计问题 \eqref{eq:incompressible-sfvp-discrete-pressure}
的误差, 我们将使用与 \eqref{eq:appendix-neumann-ritz-projection}
所定义投影类似的算子, 仍记为 $P_h$,
它从 $W^{1,\alpha}(\Omega)\cap L_0^2(\Omega)$ 映到 $Q_h$.
此外, 还需要下面关于向量场分解正则性的假设.

\MathBookSourcePage{223}
\medskip
\noindent\textbf{Hypothesis H2.}
每个向量场 $\mathbf v\in H_0^1(\Omega)^2$ 都有分解
\[
\mathbf v=\operatorname{grad}q+\operatorname{curl}\phi,
\qquad
\phi\in\Phi_s,
\quad q\in W^{1,s}(\Omega)\cap L_0^2(\Omega),
\]
并且
\begin{equation}\label{eq:incompressible-sfvp-decomposition-approximation}
|P_hq-q|_{1,\beta,\Omega}
+\inf_{\phi_h\in\Phi_h}|\phi_h-\phi|_{1,\beta,\Omega}
\leq Ch^{2/\beta}|\mathbf v|_{1,\Omega}
\qquad\forall\beta\in[2,s],
\end{equation}
其中常数 $C$ 与 $h$, $\phi$, $q$ 和 $\mathbf v$ 无关.

\begin{Theorem}\label{thm:incompressible-sfvp-pressure-error}
令 $p$ 和 $p_h$ 分别为问题
\eqref{eq:incompressible-decoupled-pressure}
和 \eqref{eq:incompressible-sfvp-discrete-pressure} 的解.
除 Theorem~\ref{thm:incompressible-sfvp-velocity-vorticity-error}
的假设外, 再假设 Hypothesis H2 成立, 且对某个实数
$\alpha\in[r,2]$ 有
$p\in W^{1,\alpha}(\Omega)$ 和
$\omega\in W^{1,\alpha}(\Omega)$.
则压力误差满足
\begin{equation}\label{eq:incompressible-sfvp-pressure-error}
\begin{split}
\|p-p_h\|_{0,\Omega}
\leq C\biggl\{&h^{2/\beta}
\inf_{q_h\in Q_h}|p-q_h|_{1,\alpha,\Omega}
+\nu\|\omega-\omega_h\|_{0,\Omega}\\
&+\nu\inf_{\theta_h\in\Theta_h}
\left[
hK_2(h)\|\omega_h-\theta_h\|_{0,\Omega}
+h^{2/\beta}|\theta_h-\omega|_{1,\alpha,\Omega}
\right]\biggr\},
\end{split}
\end{equation}
其中 $1/\alpha+1/\beta=1$.
\end{Theorem}

\begin{Proof}
由于 $p-p_h\in L_0^2(\Omega)$, 存在
$\mathbf v\in H_0^1(\Omega)^2$, 使
\[
\operatorname{div}\mathbf v=p-p_h,
\qquad
|\mathbf v|_{1,\Omega}
\leq C\|p-p_h\|_{0,\Omega}.
\]
由 Hypothesis H2,
\[
\mathbf v=\operatorname{grad}q+\operatorname{curl}\phi,
\]
并且
\[
\|p-p_h\|_{0,\Omega}^2
=-(\operatorname{grad}(p-p_h),\mathbf v)
=(\operatorname{grad}(p_h-p),\operatorname{grad}q).
\]
引入 $P_hq$, 可写成
\[
\begin{split}
\|p-p_h\|_{0,\Omega}^2
={}&(\operatorname{grad}(q_h-p),
\operatorname{grad}(q-P_hq))\\
&+(\operatorname{grad}(p_h-p),
\operatorname{grad}(P_hq))
\qquad\forall q_h\in Q_h.
\end{split}
\]
由 \eqref{eq:incompressible-decoupled-pressure}
和 \eqref{eq:incompressible-sfvp-discrete-pressure},
\[
(\operatorname{grad}(p-p_h),\operatorname{grad}(P_hq))
=\nu(\operatorname{curl}(\omega_h-\omega),
\operatorname{grad}(P_hq)).
\]
因此, 任取 $\theta_h\in\Theta_h$ 和 $\phi_h\in\Phi_h$,
并置
\[
\mathbf v_h=\operatorname{grad}(P_hq)+\operatorname{curl}\phi_h,
\]
由 \eqref{eq:incompressible-stream-vorticity-pair}
和 \eqref{eq:incompressible-sfvp-practical-problem}, 得
\[
\begin{split}
(\operatorname{grad}(p-p_h),\operatorname{grad}(P_hq))
={}&\nu(\operatorname{curl}(\omega_h-\omega),
\mathbf v_h-\operatorname{curl}\phi_h)\\
={}&\nu(\operatorname{curl}(\omega_h-\theta_h),
\mathbf v_h-\mathbf v)\\
&+\nu(\operatorname{curl}(\theta_h-\omega),
\mathbf v_h-\mathbf v)\\
&+\nu(\omega_h-\omega,\operatorname{curl}\mathbf v)
\qquad\forall\theta_h\in\Theta_h,
\ \forall\phi_h\in\Phi_h.
\end{split}
\]

\MathBookSourcePage{224}
汇总这些等式, 得
\[
\begin{split}
\|p-p_h\|_{0,\Omega}^2
={}&(\operatorname{grad}(q_h-p),
\operatorname{grad}(q-P_hq))
-\nu(\omega_h-\omega,\operatorname{curl}\mathbf v)\\
&-\nu(\operatorname{curl}(\omega_h-\theta_h),
\mathbf v_h-\mathbf v)
-\nu(\operatorname{curl}(\theta_h-\omega),
\mathbf v_h-\mathbf v)
\end{split}
\]
对所有 $q_h\in Q_h$, $\theta_h\in\Theta_h$
和 $\phi_h\in\Phi_h$ 成立.
于是 \eqref{eq:incompressible-sfvp-decomposition-approximation}
蕴含 \eqref{eq:incompressible-sfvp-pressure-error}.
\end{Proof}

最后给出空间 $\Phi_h$, $\Theta_h$, $M_h$ 和 $Q_h$ 的一个具体例子.
为简化讨论, 假设 $\Omega$ 是 $\mathbb R^2$ 中的多边形区域.
令 $\mathcal T_h$ 是由三角形和/或凸四边形组成的一族剖分,
其单元直径以 $h$ 为上界.
假设当 $h$ 趨于零时, $\mathcal T_h$ 族一致正则
(参见 Definition~\ref{def:appendix-regular-triangulation}).
为简洁起见, 证明只针对完全由三角形组成的剖分给出,
但只需很少修改即可推广到 $\mathcal T_h$ 还含四边形的情形.

给定整数 $l\geq1$, 选择下列标准有限元空间:
\begin{equation}\label{eq:incompressible-sfvp-standard-fe-spaces}
\left\{
\begin{aligned}
\Theta_h=M_h
&=\{\theta_h\in C^0(\overline\Omega);
\ \theta_h|_\kappa\in P_l
\ \forall\kappa\in\mathcal T_h\}
\subset W^{1,\infty}(\Omega),\\
\Phi_h
&=\Theta_h\cap\Phi
=\{\phi_h\in\Theta_h;
\ \phi_h|_{\Gamma_0}=0,
\ \phi_h|_{\Gamma_i}=\text{an arbitrary constant }c_i,
\ 1\leq i\leq p\},\\
Q_h
&=\{q_h\in C^0(\overline\Omega)\cap L_0^2(\Omega);
\ q_h|_\kappa\in P_k
\ \forall\kappa\in\mathcal T_h\},
\qquad k=\min(l,l-1).
\end{aligned}
\right.
\end{equation}
压力取 $l-1$ 次多项式的理由可从界
\eqref{eq:incompressible-sfvp-pressure-error} 看出.
事实上, 无论 $p_h$ 的次数多高, $p-p_h$ 的收敛阶都不可能高于
$\omega-\omega_h$ 的收敛阶.

这里 $\Phi_h\subset\Theta_h=M_h$.
算子 $P_h$ 的逼近性质由 Theorem~\ref{thm:appendix-ritz-projection-lp-error}
给出, 而 $K_r(h)$ 的估计由 Corollary~\ref{cor:appendix-global-inverse-gradient}
给出; 这两个结果都要求 $\mathcal T_h$ 一致正则:
\begin{equation}\label{eq:incompressible-sfvp-inverse-estimate}
\|\theta_h\|_{1,r,\Omega}
\leq C(r)\frac1h\|\theta_h\|_{0,\Omega}
\qquad\forall\theta_h\in\Theta_h,
\quad\forall r\leq2,
\end{equation}
其中常数 $C(r)$ 依赖于 $r$ 和 $\Omega$.
因此, 对 $\widetilde V_h$ 中的逼近误差有估计
\begin{equation}\label{eq:incompressible-sfvp-kernel-approximation}
\inf_{\mathbf v_h\in\widetilde V_h}
\|\mathbf v-\mathbf v_h\|_{\widetilde X}
\leq Ch^{k-1}
\left(|\phi|_{k+1,s,\Omega}+\|\theta\|_{k,\Omega}\right),
\quad 0\leq k\leq l,
\end{equation}
对所有
$\mathbf v=(\operatorname{curl}\phi,\theta)
\in\widetilde V\cap
(W^{k+1,s}(\Omega)\times H^k(\Omega))$ 成立.

于是只需验证 Hypothesis H1 (当 $s>2$ 时) 和 Hypothesis H2.
下面两个 Lemma 完成这一工作.
二者都假设 $\Omega$ 是凸多边形,
因为证明需要 Laplace 算子 Dirichlet 问题之解具有额外正则性.

\MathBookSourcePage{225}
\begin{Lemma}\label{lem:incompressible-sfvp-hypothesis-h1}
除上述假设外, 再假设 $\Omega$ 是凸多边形.
则对每个实数 $s>2$, 存在与 $h$ 无关的正常数 $C_s$, 使
\begin{equation}\label{eq:incompressible-sfvp-h1-bound}
|\phi_h|_{1,s,\Omega}
\leq C_s\|\theta_h\|_{0,\Omega}
\qquad
\forall\mathbf v_h=(\operatorname{curl}\phi_h,\theta_h)
\in\widetilde V_h.
\end{equation}
\end{Lemma}

\begin{Proof}
首先注意, 由于 $\Omega$ 是凸的, 其边界 $\Gamma$
只有一个连通分支, 即 $\Gamma=\Gamma_0$, 并且
$\Phi_h\subset W_0^{1,s}(\Omega)$.

其次, 由 \eqref{eq:incompressible-weak-four-field-kernel}
和 Sobolev Theorem~\ref{thm:sobolev-embedding},
$\widetilde V$ 中每个函数
$\mathbf v=(\operatorname{curl}\phi,\theta)$
都满足 \eqref{eq:incompressible-sfvp-h1-bound}.
此外,
\begin{equation}\label{eq:incompressible-sfvp-continuous-poisson-relation}
-\Delta\phi=\theta\quad\text{in }\Omega.
\end{equation}
当然, 对 $\widetilde V_h$ 中的函数该等式并不成立,
但它们满足
\[
(\operatorname{curl}\phi_h,\operatorname{curl}\mu_h)
=(\theta_h,\mu_h)
\qquad\forall\mu_h\in\Theta_h,
\]
这可解释为 \eqref{eq:incompressible-sfvp-continuous-poisson-relation}
的离散版本. 这提示我们把 $\phi_h$ 与辅助函数 $\phi(h)$ 比较,
其中 $\phi(h)$ 是齐次 Dirichlet 问题
\[
-\Delta\phi(h)=\theta_h\quad\text{in }\Omega,
\qquad
\phi(h)=0\quad\text{on }\Gamma
\]
的解.
由于 $\Omega$ 是凸的, Theorem~\ref{thm:laplace-dirichlet-regularity}
断言 $\phi(h)\in H^2(\Omega)$, 且
\begin{equation}\label{eq:incompressible-sfvp-auxiliary-h2-bound}
\|\phi(h)\|_{2,\Omega}
\leq C_1\|\theta_h\|_{0,\Omega}.
\end{equation}
此外, 容易看出 $\phi_h=P_h\phi(h)$,
其中 $P_h$ 是 $\phi(h)$ 到 $\Phi_h$ 上的 $H^1$ 投影
(参见 \eqref{eq:appendix-dirichlet-ritz-projection}).
因此 Theorem~\ref{thm:appendix-ritz-projection-lp-error} 给出
\[
|\phi(h)-\phi_h|_{1,s,\Omega}
\leq C_2\|\phi(h)\|_{1,s,\Omega}.
\]
再应用 \eqref{eq:incompressible-sfvp-auxiliary-h2-bound}
和 Sobolev Theorem~\ref{thm:sobolev-embedding}, 得
\[
|\phi_h|_{1,s,\Omega}
\leq(1+C_2)\|\phi(h)\|_{1,s,\Omega}
\leq C_s\|\theta_h\|_{0,\Omega}.
\]
\end{Proof}

\begin{Remark}\label{rem:incompressible-sfvp-elementary-h1-proof}
Lemma~\ref{lem:incompressible-sfvp-hypothesis-h1}
还有一个不调用 Theorem~\ref{thm:appendix-ritz-projection-lp-error}
的更初等证明, 后者是一个困难结果.
一方面, $\phi_h=P_h\phi(h)$ 蕴含
\[
|\phi_h-\chi_h|_{1,\Omega}
\leq|\phi(h)-\chi_h|_{1,\Omega}
\qquad\forall\chi_h\in\Phi_h,
\]
而 Lemma~\ref{lem:appendix-global-inverse-same-order} 对所有 $s>2$ 给出
\[
|\phi_h-\chi_h|_{1,s,\Omega}
\leq C(s)h^{2/s-1}|\phi_h-\chi_h|_{1,\Omega}.
\]
另一方面, 可写成
\[
|\phi_h|_{1,s,\Omega}
\leq|\phi_h-\chi_h|_{1,s,\Omega}
+|\chi_h|_{1,s,\Omega}
\qquad\forall\chi_h\in\Phi_h.
\]
汇总这三个不等式, 得
\[
|\phi_h|_{1,s,\Omega}
\leq\inf_{\chi_h\in\Phi_h}
\left\{C(s)h^{2/s-1}|\phi(h)-\chi_h|_{1,\Omega}
+|\chi_h|_{1,s,\Omega}\right\}.
\]

\MathBookSourcePage{226}
由于 $s>2$, Sobolev Theorem~\ref{thm:sobolev-embedding}
断言 $W^{1,s}(\Omega)\subset C^0(\overline\Omega)$.
因此可取 $\chi_h=I_h\phi(h)$,
其中 $I_h$ 是 Lemma~\ref{lem:appendix-lagrange-interpolation}
定义的标准有限元插值算子. 该 Lemma 给出
\[
|I_h\phi(h)|_{1,s,\Omega}
\leq C_s|\phi(h)|_{1,s,\Omega},
\qquad
|\phi(h)-I_h\phi(h)|_{1,\Omega}
\leq C_1h|\phi(h)|_{2,\Omega}.
\]
结合 \eqref{eq:incompressible-sfvp-auxiliary-h2-bound},
即得 \eqref{eq:incompressible-sfvp-h1-bound}.
\end{Remark}

\begin{Remark}\label{rem:incompressible-sfvp-nonconvex-h1}
当 $s$ 不太大时, 可以证明
\eqref{eq:incompressible-sfvp-h1-bound}
并不要求 $\Omega$ 为凸域.
注意, 上述另一种证明仅用这一假设推出 $\phi(h)\in H^2(\Omega)$,
而实际上它并不需要这么高的正则性.
例如, 若 $s=4$ (这是 Navier--Stokes 方程最常使用的指数),
只需 $\phi(h)\in H^{3/2}(\Omega)$,
因为 $H^{3/2}(\Omega)\subset W^{1,4}(\Omega)$.
但 Theorem~\ref{thm:laplace-dirichlet-regularity}
的推广 (参见 Grisvard~\cite{ref:grisvard-42}) 证明,
只要 $\Gamma$ 的各内角以 $3\pi/2$ 为上界,
$\phi(h)$ 就属于 $H^{3/2}(\Omega)$,
这允许更广泛的一类区域.
读者可在 Subsection~\ref{subsec:navier-centered-stream-vorticity} 中找到更多细节.
\end{Remark}

\begin{Lemma}\label{lem:incompressible-sfvp-hypothesis-h2}
当 $\Omega$ 是凸多边形且 $\mathcal T_h$ 一致正则时,
Hypothesis H2 对所有实数 $s>2$ 成立.
\end{Lemma}

\begin{Proof}
令 $\mathbf v\in H_0^1(\Omega)^2$.
使用 Theorem~\ref{thm:sec3-orthogonal-decomposition}
建立的正交分解
\[
\mathbf v=\operatorname{grad}q+\operatorname{curl}\phi,
\]
其中 $\phi\in\Phi$, 并且由
Remark~\ref{rem:sec3-neumann-interpretation},
$q\in H^1(\Omega)\cap L_0^2(\Omega)$ 是问题
\[
\Delta q=\operatorname{div}\mathbf v\quad\text{in }\Omega,
\qquad
\frac{\partial q}{\partial n}=0\quad\text{on }\Gamma
\]
的解.
由于 $\Omega$ 是凸的,
Theorem~\ref{thm:laplace-neumann-regularity}
断言 $q\in H^2(\Omega)\cap L_0^2(\Omega)$, 且
\[
\|q\|_{2,\Omega}
\leq C_1\|\operatorname{div}\mathbf v\|_{0,\Omega}
\leq C_1\sqrt2|\mathbf v|_{1,\Omega}.
\]
于是 $\phi\in H^2(\Omega)\cap H_0^1(\Omega)$, 且
\[
\|\phi\|_{2,\Omega}
\leq C_2\|\mathbf v\|_{1,\Omega}
\leq C_3|\mathbf v|_{1,\Omega}.
\]
由 Sobolev Theorem~\ref{thm:sobolev-embedding},
\[
H^2(\Omega)\subset W^{1+2/s,s}(\Omega).
\]
由于 $Q_h$ 中的多项式至少为一次,
\eqref{eq:appendix-ritz-projection-lp-error}
蕴含
\[
|q-P_hq|_{1,s,\Omega}
\leq C_4h^{2/s}|\mathbf v|_{1,\Omega}.
\]

\MathBookSourcePage{227}
同样, \eqref{eq:appendix-lagrange-interpolation-error-wsp}
和 Jensen 不等式 \eqref{eq:appendix-discrete-jensen}
给出
\[
|\phi-I_h\phi|_{1,s,\Omega}
\leq C_5h^{2/s}|\mathbf v|_{1,\Omega}.
\]
因此 Hypothesis H2 对所有 $s>2$ 成立.
\end{Proof}

由 Theorem~\ref{thm:incompressible-sfvp-velocity-vorticity-error},
Theorem~\ref{thm:incompressible-sfvp-pressure-error}
和上面两个 Lemma, 立即得到本小节的主要结果.

\begin{Theorem}\label{thm:incompressible-sfvp-standard-fe-error}
令 $\Omega$ 是有界开凸多边形,
$\mathcal T_h$ 是一族一致正则剖分.
令
$(\mathbf u=(\operatorname{curl}\psi,\omega=-\Delta\psi),p)$
为问题 \eqref{eq:incompressible-stream-vorticity-pair}
和 \eqref{eq:incompressible-decoupled-pressure} 的解, 并假设
\begin{equation}\label{eq:incompressible-sfvp-standard-regularity}
\psi\in W^{k+1,s}(\Omega),
\qquad
\Delta\psi\in H^k(\Omega),
\qquad
p\in H^m(\Omega)\cap L_0^2(\Omega),
\qquad
m=\max(1,k-1),
\end{equation}
其中整数 $k$ 满足 $1\leq k\leq l$.
则有误差界
\begin{equation}\label{eq:incompressible-sfvp-standard-velocity-error}
\|\mathbf u-\mathbf u_h\|_{\widetilde X}
\leq C_1\left[
h^{k-1}\left(|\psi|_{k+1,s,\Omega}+|\Delta\psi|_{k,\Omega}\right)
+h^k\|\Delta\psi\|_{k,\Omega}
\right],
\end{equation}
以及
\begin{equation}\label{eq:incompressible-sfvp-standard-pressure-error}
\|p-p_h\|_{0,\Omega}
\leq C_2\left[
h^{k-1}\left(|p|_{m,\Omega}+|\psi|_{k+1,s,\Omega}
+|\Delta\psi|_{k,\Omega}\right)
+h^k\|\Delta\psi\|_{k,\Omega}
\right],
\end{equation}
其中常数 $C_1>0$ 和 $C_2>0$ 与 $h$, $\psi$ 和 $p$ 无关.
\end{Theorem}

\begin{Remark}\label{rem:incompressible-sfvp-convexity-and-suboptimality}
Theorem~\ref{thm:incompressible-sfvp-standard-fe-error}
在三个地方使用了 $\Omega$ 的凸性:
验证 Hypothesis H1 和 Hypothesis H2,
以及估计 $\|\omega-P_h\omega\|_{0,\Omega}$.
在许多情形下, Hypothesis H1 不要求 $\Omega$ 为凸域
(参见 Remark~\ref{rem:incompressible-sfvp-nonconvex-h1}).
但 $P_h\omega$ 的 $L^2$ 估计来自对偶论证,
而 Hypothesis H2 也要求 $\Omega$ 为凸域, 因为 $\Gamma$ 有角点.
就 \eqref{eq:incompressible-sfvp-standard-velocity-error} 而言,
总可以用
$\inf_{\theta_h\in\Theta_h}\|\omega-\theta_h\|_{1,r,\Omega}$
代替 $\|\omega-P_h\omega\|_{0,\Omega}$
(参见 \eqref{eq:incompressible-velocity-abstract-error}),
并且若 $\omega$ 足够光滑, 即使 $\Omega$ 不凸,
也能为 $\|\mathbf u-\mathbf u_h\|_{\widetilde X}$
导出同阶估计.

但遗憾的是, 我们对
\eqref{eq:incompressible-sfvp-standard-pressure-error}
的证明确实要求 $\Omega$ 为凸域, 而且看不到去掉这一假设的可能性.

由于存在因子 $K(h)$,
Theorem~\ref{thm:incompressible-sfvp-standard-fe-error}
的估计并非最优. 特别地, 当 $l=1$, 即使用分片线性元时,
\eqref{eq:incompressible-sfvp-standard-velocity-error}
和 \eqref{eq:incompressible-sfvp-standard-pressure-error}
并不蕴含收敛. 如前所述, Subsection~\ref{subsec:incompressible-sfvp-mesh-dependent}
将处理这一困难.
\end{Remark}
